\documentclass[14pt]{beamer}
\usepackage{pepslides}
\title{Spherical Earth structure}
\lecturenumber{18}
\hypersetup{pdftitle={Lecture 18: Spherical Earth structure}}

% Ray from hypocentre to seismometer and the epicentral angle (used on two slides).
\newcommand{\raygeometry}{%
\begin{tikzpicture}[x=1.05cm,y=1.05cm,line width=0.9pt]
\coordinate (O) at (0,0);
\coordinate (E) at (128:5); \coordinate (S) at (52:5); \coordinate (F) at (128:4.3);
\draw (142:5) arc (142:38:5);
\draw[dashed] (O) -- (E); \draw[dashed] (O) -- (S);
\draw[line width=0.6pt] (52:0.7) arc (52:128:0.7);
\node at (90:1.15) {$\Delta$};
\draw[postaction={decorate},decoration={markings,mark=at position 0.55 with {\arrow{Stealth[length=3mm]}}}]
  (F) to[bend right=16] (S);
\fill (F) circle (0.08);
\node[anchor=west] at ([shift={(0.12,0.12)}]F) {Hypocentre};
\draw[line width=0.6pt] (E) -- ++(0,1.1) node[above] {Epicentre};
\draw[line width=0.6pt] (S) -- ++(0,1.1) node[above] {Seismometer};
\node[below,align=center] at (0,-0.15) {Centre of\\the Earth};
\end{tikzpicture}}

\begin{document}

\titleslide

\begin{frame}[t]
\heading{From last time\ldots}
Ray Hamiltonian defined by:
\begin{eqn}
H_{k}(\mathbf{x},\mathbf{p})=\frac{1}{2}\|\mathbf{p}\|^{2}c_{k}(\mathbf{x},\hat{\mathbf{p}})^{2}.
\end{eqn}
Eikonal equation:
\begin{eqn}
H_{k}[\mathbf{x},\mathbf{p}(\mathbf{x})]=\frac{1}{2}, \quad \mathbf{p}(\mathbf{x})=(\nabla T)(\mathbf{x}).
\end{eqn}
Solved through integration of Hamilton's equations:
\begin{eqn}
\frac{\ddns x_{i}}{\ddns\sigma}=\frac{\partial H_{k}}{\partial p_{i}}, \quad \frac{\ddns p_{i}}{\ddns\sigma}=-\frac{\partial H_{k}}{\partial x_{i}}.
\end{eqn}
\end{frame}

\sectionslide{Earthquake location}

\begin{frame}
\centering
\includegraphics[height=0.42\paperheight]{L6F1.pdf}

\raggedright
The \textbf{hypocentre} is where the rupture initiates. The epicentre is the point on the surface directly above it.

It is the origin of the \textbf{first waves} recorded on seismometers.
\end{frame}

\begin{frame}
\centering
\includegraphics[height=0.58\paperheight]{L6F2.png}

\raggedright
The \textbf{arrival time} is the actual time (e.g.\ noon on Wednesday) at which a seismic phase arrives at a seismometer.
Measured from \textbf{high-frequency} recordings to allow for precise measurements.
\end{frame}

\begin{frame}
Arrival time at the $i$th seismometer for a given phase takes the form:
\begin{eqn}
t_{i}=t_{h}+T(\mathbf{x}_{i},\mathbf{x}_{h})+e_{i}.
\end{eqn}
Travel time can be calculated efficiently using ray theory. But to do this, we need to know the Earth's velocity structure.

Forward problem depends on just four model parameters:
\begin{itemize}
\item Three components of the hypocentre position, $\mathbf{x}_{h}$
\item The hypocentral time, $t_{h}$
\end{itemize}
\end{frame}

\begin{frame}
Misfit between data and proposed model parameters can be quantified:
\begin{eqn}
J(\mathbf{x}_h, t_h) = \sum_{i=1}^{N_s} \frac{[t_h + T(\mathbf{x}_i, \mathbf{x}_h) - t_i]^2}{\varsigma_i^2},
\end{eqn}
where $\varsigma_{i}^{2}$ is the variance of the measurement error $e_{i}$.

Optimise $J$ with respect to the four model parameters to find a solution. This can be done iteratively using
gradient-based methods or through a brute-force grid search.

Uncertainties can be estimated by considering a range of model parameters that fit the data sufficiently well.

Accuracy ultimately limited by the \textbf{dominant wavelength} of observed waves.
\end{frame}

\begin{frame}
\centering
\includegraphics[height=0.62\paperheight]{pavlis2000_pamir_map.jpg}\hspace{2em}%
\includegraphics[height=0.55\paperheight]{pavlis2000_pamir_sections.jpg}

\raggedright
Application of hypocentral locations to test competing tectonic models for the formation of the Pamir--Hindu Kush mountains.
\credit{Pavlis \& Das (2000)}
\end{frame}

\sectionslide{Rays in spherical earth models}

\begin{frame}
We focus on an isotropic model, and look at P-waves for definiteness.

Eikonal equation takes the form:
\begin{eqn}
\begin{gathered}
H(\mathbf{x},\mathbf{p})=\frac{1}{2},\\[\eqgap]
H(\mathbf{x},\mathbf{p})=\frac{1}{2}\alpha(\mathbf{x})^{2}\|\mathbf{p}\|^{2}.
\end{gathered}
\end{eqn}
We assume further that the P-wave velocity \textbf{varies only with radius}.
\end{frame}

\begin{frame}
General form of Hamilton's equations:
\begin{eqn}
\frac{\ddns \mathbf{x}}{\ddns \sigma}=\frac{\partial H}{\partial \mathbf{p}},\quad \frac{\ddns \mathbf{p}}{\ddns \sigma}=-\frac{\partial H}{\partial \mathbf{x}},
\end{eqn}
with the travel time along a ray given by $T[\mathbf{x}(\sigma)]=\sigma$.

For the spherically symmetric P-wave problem, they become:
\begin{eqn}
\frac{\ddns \mathbf{x}}{\ddns \sigma}=\alpha\hat{\mathbf{p}},\quad \frac{\ddns \mathbf{p}}{\ddns \sigma}=-\frac{1}{\alpha}\frac{\ddns \alpha}{\ddns r}\hat{\mathbf{x}}.
\end{eqn}
Here hats denote unit vectors.
\end{frame}

\begin{frame}
We define the vector:
\begin{eqn}
\mathbf{q}=\mathbf{x}\times \mathbf{p}.
\end{eqn}
Using Hamilton's equations, we find that:
\begin{eqn}
\frac{\ddns }{\ddns \sigma}(\mathbf{x}\times \mathbf{p})=\frac{\ddns \mathbf{x}}{\ddns \sigma}\times \mathbf{p}+\mathbf{x}\times\frac{\ddns \mathbf{p}}{\ddns \sigma} =
\alpha\hat{\mathbf{p}}\times \mathbf{p}-\frac{1}{\alpha}\frac{\ddns \alpha}{\ddns r}\mathbf{x}\times\hat{\mathbf{x}}=\mathbf{0}.
\end{eqn}
The form of this conserved quantity follows from the rotational symmetry via Noether's theorem (but you don't have to know this!)
\end{frame}

\begin{frame}
Conservation of $\mathbf{q}$ implies that the ray stays in the plane spanned by the initial values of $\mathbf{x}$ and $\mathbf{p}$.

Dynamics \textbf{reduced to a plane} -- cf.\ the two-body problem.

Magnitude of the conserved quantity is:
\begin{eqn}
q=\frac{r \sin \theta}{\alpha},
\end{eqn}
where $\theta$ is the angle between $\mathbf{x}$ and $\mathbf{p}$.

This is \textbf{Snell's law} expressed in spherical geometry.
\end{frame}

\begin{frame}
Along the ray, we can write:
\begin{eqn}
\hat{\mathbf{p}}=\cos[\theta(\mathbf{x})]\hat{\mathbf{x}}+\sin[\theta(\mathbf{x})]\hat{\bm{\Delta}}(\mathbf{x}),
\end{eqn}
where $\Delta$ is the epicentral angle.

\vspace{0.5em}
\centering
\raygeometry
\end{frame}

\begin{frame}
From the conserved quantity, we can write:
\begin{eqn}
\sin \theta=\frac{\alpha q}{r},\quad \cos \theta=\pm\sqrt{1-\frac{\alpha^{2}q^{2}}{r^{2}}},
\end{eqn}
showing that the ray's direction is fixed by its radius. Plus or minus sign in the cosine term corresponds to up- or down-going rays.

Using these results in Hamilton's equations, we arrive at (for a down-going ray):
\begin{eqn}
\frac{\ddns \mathbf{x}}{\ddns \sigma}=-\alpha\sqrt{1-\frac{\alpha^{2}q^{2}}{r^{2}}}\hat{\mathbf{x}}+\frac{\alpha^{2}q}{r}\hat{\bm{\Delta}}.
\end{eqn}
The slowness vector has now been eliminated from the problem.
\end{frame}

\begin{frame}
Using the reduced ODE, we find that:
\begin{eqn}
\begin{gathered}
\frac{\ddns }{\ddns \sigma}r[\mathbf{x}(\sigma)]=\frac{\partial r}{\partial x_{i}}\frac{\ddns x_{i}}{\ddns\sigma}=-\alpha\sqrt{1-\frac{\alpha^{2}q^{2}}{r^{2}}},\\[\eqgap]
\frac{\ddns }{\ddns \sigma}\Delta[\mathbf{x}(\sigma)]=\frac{\partial\Delta}{\partial x_{i}}\frac{\ddns x_{i}}{\ddns \sigma}=\frac{\alpha^{2}q}{r^{2}}.
\end{gathered}
\end{eqn}
And hence, we have:
\begin{eqn}
\frac{\ddns \Delta}{\ddns r}=\frac{\ddns \Delta}{\ddns\sigma}\left(\frac{\ddns r}{\ddns \sigma}\right)^{-1}=\frac{-\alpha q}{r\sqrt{r^{2}-\alpha^{2}q^{2}}}.
\end{eqn}
\end{frame}

\begin{frame}
In a similar manner, using $\ddns T/\ddns\sigma=1$, we find:
\begin{eqn}
\frac{\ddns T}{\ddns r}=\frac{\ddns T}{\ddns \sigma}\left(\frac{\ddns r}{\ddns \sigma}\right)^{-1}=\frac{-r}{\alpha\sqrt{r^{2}-\alpha^{2}q^{2}}}.
\end{eqn}
Valid only where $\Delta$ and $T$ are single-valued functions of $r$.

\vspace{0.5em}
\centering
\raygeometry
\end{frame}

\begin{frame}
Assuming velocity decreases monotonically with radius, the ray reaches a \textbf{turning point} where it is travelling horizontally:
\begin{eqn}
q=\frac{r_{t}}{\alpha(r_{t})}.
\end{eqn}
Afterwards, it returns to the surface, $r=b$, along a symmetric path.

Angle and travel time expressed through the integrals:
\begin{eqn}
\Delta=\int_{r_{t}}^{b}\frac{2\alpha(r)q}{r\sqrt{r^{2}-\alpha(r)^{2}q^{2}}}\dd r,\quad
T=\int_{r_{t}}^{b}\frac{2r}{\alpha(r)\sqrt{r^{2}-\alpha(r)^{2}q^{2}}}\dd r.
\end{eqn}
These expressions could be evaluated in the 1920s and 1930s using \textbf{mechanical calculating machines}.
\end{frame}

\sectionslide{Spherical velocity structure}

\begin{frame}
\centering
\includegraphics[height=0.68\paperheight]{iris_global_stack.jpg}\hspace{2em}%
\includegraphics[height=0.68\paperheight]{L6F4.pdf}

Global travel time observations (left), and travel times calculated by ray tracing in PREM (right).
\credit{IRIS}
\end{frame}

\begin{frame}
\centering
\includegraphics[width=0.9\textwidth]{rays_pkkkkp.pdf}

Ray theoretical predictions for the phase PKKKKP.
\end{frame}

\begin{frame}
\centering
\includegraphics[height=0.62\paperheight]{seismology_pioneers.jpg}

Oldham (1906): existence of the core. Gutenberg (1914): depth of the core--mantle boundary.

Jeffreys: a fluid core, from solid-Earth tides. Lehmann: a solid inner core.
\end{frame}

\begin{frame}
\begin{columns}[c]
\column{0.68\textwidth}
\centering
\includegraphics[height=0.62\paperheight]{L6F5.pdf}

\textbf{P}reliminary \textbf{R}eference \textbf{E}arth \textbf{M}odel (PREM) of Dziewonski \& Anderson (1981)
\column{0.3\textwidth}
\centering
\includegraphics[height=0.26\paperheight]{dziewonski_photo.jpg}

Adam Dziewonski (1936--2016)

\includegraphics[height=0.26\paperheight]{anderson_photo.jpg}

Don Anderson (1933--2014)
\end{columns}
\end{frame}

\begin{frame}
The P-wave and S-wave speeds are given by:
\begin{eqn}
\alpha=\sqrt{\frac{\kappa+\frac{4}{3}\mu}{\rho}},\quad \beta=\sqrt{\frac{\mu}{\rho}}.
\end{eqn}
With depth, compression increases both the elastic moduli and the density. The former effect dominates, and so wave speeds typically increase.

Sharp \textbf{discontinuities} are identified principally from reflected waves, and are due to changes in composition or phase.
\end{frame}

\begin{frame}
\centering
\includegraphics[height=0.7\paperheight]{jb_prem.pdf}

\raggedright
The model of Jeffreys and Bullen from 1940 (dashed) against PREM from 1981 (solid).
\end{frame}

\begin{frame}
\centering
\includegraphics[height=0.66\paperheight]{L6F3.pdf}\hspace{2em}%
\includegraphics[height=0.66\paperheight]{L6F4.pdf}

Existence of the Earth's core: the drop in P-wave speed at the core--mantle boundary leaves a P-wave \textbf{shadow zone}.
\end{frame}

\begin{frame}
The most notable discontinuities are:
\begin{itemize}
\item The \textbf{core--mantle boundary} (CMB): rocky mantle to fluid iron outer core
\item The \textbf{inner core boundary} (ICB): fluid to solid phase change within the iron core
\end{itemize}
Fluid outer core: no S-waves ($\mu=0$), solid-Earth tides, and the geodynamo.

Solid inner core: free oscillation periods, and (more debated) inner core S-waves.
\end{frame}

\begin{frame}
\centering
\includegraphics[width=0.9\textwidth]{rays_pcp_pcs.pdf}

Constraints on boundary locations come largely from reflected phases.
\end{frame}

\begin{frame}
\centering
\includegraphics[height=0.54\paperheight]{L6F5.pdf}\hspace{2em}%
\includegraphics[height=0.54\paperheight]{l18_mantle_mineralogy.jpg}

\raggedright
Jumps in velocity at around $410\,\mathrm{km}$ and $670\,\mathrm{km}$ depth are believed to represent \textbf{solid-state phase changes}
in silicate minerals -- you don't need to know any details!

The region between them is the \textbf{transition zone}, separating the upper and lower mantle.
\end{frame}

\begin{frame}
\centering
\includegraphics[width=0.9\textwidth]{rays_pdiff_sdiff.pdf}

There are also core-diffracted phases that are not predicted by ray theory.
\end{frame}

\begin{frame}
\centering
\includegraphics[height=0.72\paperheight]{hosseini2015_pdiff.jpg}

Observations and synthetics for core-diffracted P-waves (normalised amplitudes).
\credit{Hosseini \& Sigloch (2015)}
\end{frame}

\begin{frame}[t]
\heading{What you need to know and be able to do}
\begin{enumerate}
\item[(i)] Know what the hypocentre of an earthquake is, and how the hypocentral location and time can be determined
using observations of seismic waves.
\item[(ii)] Understand how the Hamiltonian ray tracing equations within a spherically symmetric earth model can be
reduced to the calculation of appropriate radial integrals. In the second problem set you will see how such travel
time integrals can be inverted for the velocity structure.
\item[(iii)] Understand in outline the observations and methods that were used to constrain the Earth's spherically
symmetric velocity structure. You should also know the basic properties of this spherically symmetric velocity
structure, and the physical interpretation of its major features.
\end{enumerate}
\end{frame}

\end{document}
